\begin{exercice*}
    Précise si les affirmations suivantes sont vraies ou fausses.
    \begin{enumerate}
        \item La probabilité de gagner à un jeu peut être égale à \num{1,2}.
        \item Sachant que la probabilité de gagner à un jeu est égale à \num{0,6}, la probabilité de perdre est égale à \num{1,6}.
        \item Sachant que la probabilité de perdre à un jeu est égale à \num{0,3}, la probabilité de gagner est égale à \num{0,7}.
        \item Si la probabilité d'un événement est égale à \num{0,25}, la probabilité de l'événement contraire est égale à \num{0,85}.
        \item La somme des probabilités d'un événement et de son événement contraire est égale à 1.
        \item La probabilité d'un événement est un nombre compris entre 0 et 1.
    \end{enumerate}
\end{exercice*}
\begin{corrige}
    %\setcounter{partie}{0} % Pour s'assurer que le compteur de \partie est à zéro dans les corrigés
    %\phantom{rrr}    
    \begin{enumerate}
        \item Faux. La probabilité d'un événement est un nombre compris entre 0 et 1.
        \item Faux. La probabilité de perdre est égale à $1 - \num{0,6} = \num{0,4}$.
        \item Vrai. La probabilité de gagner est égale à $1 - \num{0,3} = \num{0,7}$.
        \item Faux. La probabilité de l'événement contraire est égale à $1 - \num{0,25} = \num{0,75}$.
        \item Vrai. La somme des probabilités d'un événement et de son événement contraire est égale à $1$.
        \item Vrai. La probabilité d'un événement est un nombre compris entre $0$ et $1$.
    \end{enumerate}
\end{corrige}

